\section{The inequality and its geometric quantity}
\label{sec:introduction}
The Penrose inequality asks how much mass is required to surround a
trapped region. Its origin is a test of cosmic censorship: a collapsing
configuration should settle to a black hole, radiation should carry
energy away, and the horizon area should not decrease. Penrose's
null-shell argument made this comparison concrete \cite{Penrose1973}.
An initial-data theorem must express it using only a Riemannian metric,
a second fundamental form and a specified enclosing area.

The choice of area is part of the problem. A surface outside an apparent
horizon can have smaller area than the horizon itself. Ben-Dov gave
spherically symmetric examples satisfying the dominant energy condition
for which the bound using the apparent horizon's own area fails
\cite{BenDov2004}. Carrasco and Mars found a different obstruction
for generalized apparent horizons: their counterexamples already have
an outer area-minimizing generalized horizon \cite{CarrascoMars2010}.
Thus a mass--area statement must specify both its trapping condition
and its enclosing class. We use instead the infimum of the areas of all full cuts
separating the boundary from infinity. Every component counts, including
any part coincident with the boundary. The mass is the Lorentz-invariant
length of the ADM energy-momentum vector. These choices lead to the
following precise formulation.

\subsection{Data, area and statements}
\begin{definition}[Initial-data exterior]\label{def:data}
Let $n\geq3$ and let $\Omega$ be a smooth connected orientable
$n$-manifold with nonempty compact smooth boundary $S$. The metric $g$
and symmetric covariant two-tensor $K$ are smooth up to $S$, and $g$ is
complete with $S$ included. Outside a compact set, $\Omega$ is a single
coordinate end $\{x\in\R^n:|x|>R_0\}$. Thus there are no additional
ends hidden in the compact-complement condition.

Write $R_g$ for scalar curvature and $\nabla$ for the connection of $g$.
Write $\tau=\tr_gK$. Our constraint normalization is
\begin{equation}\label{eq:setting:constraints}
 2\mu=R_g+(\tr_gK)^2-|K|_g^2,\qquad
 J_i=\nabla^j\bigl(K_{ij}-(\tr_gK)g_{ij}\bigr).
\end{equation}
The dominant energy condition is $\mu\geq|J|_g$. We require
$\mu,|J|_g\in L^1(\Omega,dV_g)$. Along $S$ the unit normal $\nu$
points into $\Omega$, toward the end. Its future expansion is
\[
 \theta_+=H_S+\tr_SK,\qquad
 H_S=\divg_S\nu,\qquad
 \tr_SK=(g^{ij}-\nu^i\nu^j)K_{ij}.
\]
The boundary is weakly future trapped when $\theta_+\leq0$ on every
component. No condition is imposed on the other null expansion.

For coordinate tensors, $T=O_j(r^{-a})$ means
$|\partial^\alpha T|\leq C_\alpha r^{-a-|\alpha|}$ for
$|\alpha|\leq j$, where $r=|x|$. The decay orders will be specified
in each theorem. Set $k=n-1$ and $\omega=\omega_{n-1}=|\mathbb S^{n-1}|$.
The Laplacian is $\Delta=\divg\nabla$ and symmetrization includes a factor $1/2$. We require the
finite ADM limits
\begin{align}
 E&=\frac{1}{2(n-1)\omega_{n-1}}\lim_{R\to\infty}
  \int_{|x|=R}(\partial_jg_{ij}-\partial_i g_{jj})
       n_\delta^i\,dA_\delta,\label{eq:setting:energy}\\
 P_i&=\frac{1}{(n-1)\omega_{n-1}}\lim_{R\to\infty}
  \int_{|x|=R}\bigl(K_{ij}-(\tr_gK)g_{ij}\bigr)
       n_\delta^j\,dA_\delta.\label{eq:setting:momentum}
\end{align}
Here $n_\delta$ and $dA_\delta$ are Euclidean. If $E>|P|_\delta$, define
$m=(E^2-|P|_\delta^2)^{1/2}$.

\end{definition}

\begin{definition}[Full enclosing cuts]\label{def:fullcuts}
A \emph{full enclosing cut} is the entire intrinsic manifold boundary
$\Gamma=\partial D$ of a connected smooth codimension-zero
submanifold-with-boundary $D\subset\Omega$. We require $D$ to be
closed in $\Omega$, its manifold interior to lie in
$\operatorname{int}\Omega$, and $D$ to contain the whole sufficiently
distant end. Its full boundary must be compact, smooth, embedded and
two-sided. In particular, $D=\Omega$ is admissible and has boundary $S$.
Define
\begin{equation}\label{eq:setting:minarea}
 A_{\min,g}(S)=\inf_D\Area_g(\partial D).
\end{equation}
The use of intrinsic boundary ensures that coincident portions of $S$
are counted. A cut need not be connected, trapped or minimal. The
infimum need not be attained.
\end{definition}

\begin{theorem}[The enclosing-area Penrose inequality]\label{thm:main}
Let $(\Omega^n,g,K)$ be the smooth one-ended exterior specified above,
with integrable constraint densities, finite ADM limits, dominant
energy condition and weakly future trapped boundary. Suppose
\[
 g-\delta=O_6(r^{-q}),\qquad K=O_5(r^{-1-q}),\qquad
 \frac{n-2}{2}<q<n-2,
\]
and assume $E>|P|_\delta$ and $A_{\min,g}(S)>0$. Then
\begin{equation}\label{eq:setting:main}
 m\geq\frac12\left(\frac{A_{\min,g}(S)}{\omega_{n-1}}
                 \right)^{\frac{n-2}{n-1}}.
\end{equation}
The coefficient and exponent are sharp in each dimension $n\geq3$.
\end{theorem}

Thus the enclosing-area form of the spacetime Penrose conjecture holds
in this class. Neither outermostness nor outer area minimization of $S$
is assumed. The second fundamental form is arbitrary, and no spin,
vacuum or interior-topology hypothesis is imposed. The time-symmetric
Schwarzschild--Tangherlini exterior of mass $m>0$ has
$A_{\min,g}(S)=\omega_{n-1}(2m)^{(n-1)/(n-2)}$ and attains equality.

\begin{corollary}[Automatic future timelikeness in the strong-decay class]
\label{cor:automatic-timelikeness}
Assume all hypotheses of Theorem~\ref{thm:main} except the inequality
$E>|P|_\delta$. In particular, retain the one-ended exterior, the
$O_6/O_5$ decay with $(n-2)/2<q<n-2$, and $A_{\min,g}(S)>0$.
Then $E>|P|_\delta$, and Equation~\eqref{eq:setting:main} holds.
\end{corollary}

The proof follows the numerical deduction in Section~\ref{sec:comparison}.

\begin{theorem}[Weaker decay in dimensions three and four]\label{thm:low}
Let $n=3$ or $n=4$. Let $(\Omega^n,g,K)$ be the smooth one-ended
exterior specified above, with integrable constraint densities, finite
ADM limits, dominant energy condition and weakly future trapped
boundary. Assume
\[
 g-\delta=O_2(r^{-q}),\qquad K=O_1(r^{-1-q}),\qquad q>\frac{n-2}{2}.
\]
Then $A_{\min,g}(S)>0$, $E>|P|_\delta$, and \eqref{eq:setting:main} holds.
\end{theorem}

Here positive enclosing area and timelikeness are both conclusions, and
no extra decay is required of $\tr_gK$. The two derivative regimes concern
the original data. After preparation, both lead to ends with estimates at every fixed derivative
order, allowing one common deformation theorem to prove the numerical
bound.

There are also designated-end formulations for finitely many ends,
with different area classes in dimensions three and four. In the
three-dimensional complete-data formulation the chosen outer region
may retain other ends; in four dimensions its connected outer domain
excludes their distant tails. Sections~\ref{sec:weak3} and~\ref{sec:weak4} give their
precise definitions and proofs. Their reduction to the one-ended inequality compares
the two enclosing infima in the needed direction; it does not identify
them. The one-ended statements above isolate the common numerical
construction. Equality for the original data requires additional
horizon assumptions and a separate variational argument.

\subsection{From time symmetry to spacetime data}
When $K=0$, the dominant energy condition becomes $R_g\geq0$ and a
marginal boundary is minimal. Geroch introduced the connected-surface
Hawking-mass monotonicity underlying smooth inverse mean curvature flow
\cite{Geroch1973}. Jang and Wald related this method to the Penrose
bound under the assumption that a classical flow runs from the inner
boundary to asymptotically round spheres \cite{JangWald1977};
Huisken and Ilmanen recount these precursors in
\cite[p.~359]{HuiskenIlmanen2001}. Two different flow methods then
established the three-dimensional Riemannian inequality. Huisken and Ilmanen's weak
inverse mean curvature flow allows jumps while preserving the Hawking
mass comparison; its bound applies to each connected component of an
outermost minimal boundary \cite{HuiskenIlmanen2001}. Bray's conformal flow changes
the metric, preserving the horizon area and making mass nonincreasing, and proves
the sharp bound for the total area of a possibly disconnected horizon
\cite{Bray2001}. Bray and Lee extended that method through dimension seven
\cite{BrayLee2009}. Their numerical theorem requires no spin assumption;
the additional spin hypothesis in the published statement belongs to
its equality conclusion.

Above dimension seven, minimizing frontiers may be singular. Bi and
Zhu extend conformal flow to this setting. Their revised theorem treats
almost-minimizing Caccioppoli boundaries with minimal regular part,
outermost minimizing enclosure and a metric extending with the stated
regularity across the boundary \cite[Theorem 1.2]{BiZhu2026v2}. Their work
supplies the ancestry of the separation, mass--capacity and convergence
arguments used for the Riemannian endpoint here. The exact version we
need allows all components of a locally minimizing and outer-minimizing
frontier; its full proof is part of the Riemannian argument described
below.

Nonzero $K$ introduces a different difficulty: neither scalar curvature
nor the boundary mean curvature has the required Riemannian sign.
Jang introduced a graph equation for extending Geroch's positive-energy
argument beyond time symmetry \cite{Jang1978}. Schoen and Yau's
minimal-surface argument established the three-dimensional Riemannian
positive-mass theorem \cite{SchoenYau1979}. Their spacetime proof combined
Jang's equation with conformal deformation and a treatment of graph
blow-up to prove positive energy for general three-dimensional initial
data \cite{SchoenYau1981}. The Penrose problem asks more of the deformation:
a lower bound for enclosing area must survive along with the mass
comparison. Malec and \'O Murchadha explained the obstruction to obtaining
both controls from the direct classical-Jang and conformal-deformation
scheme, even in spherical symmetry \cite{Malec}.

In spherical symmetry, Malec and \'O Murchadha proved the inequality
for maximal data, and Hayward obtained it from monotonicity of the
Misner--Sharp energy in untrapped regions without that maximality
restriction \cite{MalecOMurchadha1994,Hayward1996}.
Bray and Khuri introduced a warped graph and a generalized Schoen--Yau
scalar-curvature identity \cite{BrayKhuri2010,BrayKhuri2011}. On solutions
of the generalized Jang equation, the identity separates the
nonnegative dominant-energy contribution and square terms from a
divergence. It does not by itself solve the equations needed to control
that divergence. They gave a new proof in spherical symmetry and
proposed coupled systems for the general problem. Han and Khuri
established existence and boundary blow-up for prescribed warping
factors with suitable boundary and asymptotic behavior \cite{HanKhuri}. Jaracz
later constructed spherical data for which one particular coupling,
generalized Jang with zero divergence, has no smooth radial solution
with the required positive warp and asymptotics \cite{Jaracz2023}. That
obstruction concerns the specified coupling and radial solution class.

Sharp spacetime results are also known beyond the spherical
three-dimensional case. Bryden, Khuri and Sormani proved the sharp
inequality and Schwarzschild rigidity in spherical symmetry in every
spatial dimension \cite{BrydenKhuriSormani2021}. Khuri and Kunduri proved a sharp spacetime
inequality for $\mathrm{SU}(\ell+1)$-invariant data of dimension $2(\ell+1)$;
their result for the full cohomogeneity-one class is Riemannian
\cite{KhuriKunduri2025}.

For nonsymmetric three-dimensional asymptotically flat data, Allen,
Bryden, Kazaras and Khuri proved a universal invariant-mass bound with
a suboptimal constant under their asymptotic hypotheses, including extra
decay of $\tr_gK$ \cite{AllenBrydenKazarasKhuri2025}. For each end, their statement uses
at least one outermost collection of future or past marginal components
and its least enclosing area. Theorems~\ref{thm:main} and~\ref{thm:low}
give the sharp coefficient under the original-data hypotheses above,
using full enclosing cuts for any weakly future trapped boundary.

Other nonsymmetric results concern different hypotheses and area
quantities. Ellithy's inequality bounds the total area of outermost future
apparent-horizon sections under a quasi-final-state condition on the
future development, negative ingoing expansion, and a piecewise smooth
horizon tube with finitely many area-nondecreasing jumps
\cite[Theorem~4.12]{Ellithy2026}. Dong's pure-trace theorem gives the
bound for each connected generalized-horizon component
\cite[Theorem~1.2]{Dong2026Sigma}; his two-convex theorem assumes a
connected outermost past apparent horizon
\cite[Theorem~1.2]{Dong2026TwoConvex}. In the latter class the stated
asymptotic restrictions force zero ADM linear momentum, as noted in
\cite[Remark~1.3]{Dong2026TwoConvex}. These area and tensor assumptions
are distinct from the full enclosing cuts and arbitrary second fundamental
forms in our statements. Mars surveys the underlying alternatives in
formulating the initial-data Penrose problem \cite{Mars2009}.

\subsection{How the proof controls curvature, mass and area}
The numerical argument first prepares the end, then performs one filled
graph and conformal deformation. The preparation replaces the distant
end by a vacuum reference slice with matching charges, bends that slice
to rest, and repairs the constraints while comparing every full enclosing
cut. Its output has compactly supported $K$, strict dominant energy
condition and strict boundary trapping. In dimensions three and four
the weak-decay preparation provides the same differentiated output.
The three-dimensional enclosing infima converge to the original one;
in four dimensions the available lower comparison is sufficient. The
finite-stage data, rather than a formal limiting object, enter the
common deformation theorem.

Fix a prepared exterior $(\Omega,g_0,K_0)$, with energy $E_0$ and
enclosing area $A_0=A_{\min,g_0}(S)$. For each small $\varepsilon>0$,
the deformation theorem constructs a family of smooth metrics indexed
by a large parameter $N$. Each metric has a minimizing frontier lying in the original exterior
and separating its inner boundary from infinity, with area at least
$e^{-(n-1)\varepsilon}A_0$, nonnegative scalar
curvature outside that frontier, and energy at most $E_0$
plus an error tending to zero as $N\to\infty$. The Riemannian numerical
inequality therefore gives
\[
 E_0\geq
 \frac12\left(\frac{e^{-(n-1)\varepsilon}A_0}
 {\omega_{n-1}}\right)^{\frac{n-2}{n-1}}
\]
after $N\to\infty$. Letting $\varepsilon\downarrow0$ proves the
prepared inequality. Only then do we remove the end preparation.
No metric at $N=\infty$ is needed.

The construction must explain all three estimates together. We fill
the inner boundary smoothly and solve for a graph height and a conformal
factor on the resulting manifold without boundary. The conformal
factor has a lower bound, protecting the area of every enclosing cut.
The curvature identity separates the dominant energy term, coercive
squares and a divergence. In contrast to the zero-divergence coupling,
the divergence has a specified nonzero source term; its integral controls
the change of mass. Compared with the Schoen--Yau and Bray--Khuri
identities \cite{SchoenYau1981,BrayKhuri2011}, the trace equation and the
divergence source are modified to enforce these simultaneous controls.

Global solvability is the analytic part of this construction. Height
comparison and a barrier argument establish the lower conformal bound.
Testing the scalar divergence equation and applying a Sobolev inequality
on the graph, based on Michael--Simon \cite{MichaelSimon1973}, yield the
complementary upper estimate before uniform ellipticity is available.
Together these bounds control graph distortion and make the equations
uniformly elliptic for each fixed $N$.
Separate scalar regularity estimates treat the trace equation, whose
distinguished direction is the solution's gradient, and the divergence
equation, which has quadratic gradient growth. They supply the control
needed for continuation and exhaustion. These steps prove solvability
of this particular coupled system; prescribed-warp existence for another
system is not used as a substitute.

The filling removes inner boundary flux, but a minimizing enclosure
might enter the filled region. We prevent this by a further conformal
change, large deep inside the filling and small near infinity. Two-sided
density estimates force the minimizing frontier away from the interior
obstacle and into the band where scalar curvature has the correct sign.
This is why scalar-curvature improvement alone would not finish the
proof: the location and full area of the enclosure must be controlled.

For dimensions three through seven, Bray--Lee's numerical theorem
supplies the Riemannian endpoint. In higher dimensions the detached
frontier can have a singular set of dimension at most $n-8$. The required
endpoint is the minimizing-frontier theorem, including its conformal
flow, separation, smoothing and limiting arguments. We use the complete companion argument of \cite{RiemannianCompanion};
Section~\ref{sec:comparison} quotes the precise mathematical input. Its numerical application uses the stated
second-order asymptotic decay and weak-boundary hypotheses. Only the
numerical bound is needed; no Riemannian equality theorem or all-orders
expansion of the input metric at infinity is used. The endpoint thus
requires control of both the singular frontier and the asymptotically
flat end.

Section~\ref{sec:comparison} first states the prepared geometric output
and proves the numerical comparison. Section~\ref{sec:end} prepares the
strongly decaying original end. Sections~\ref{sec:scalar} and
\ref{sec:elliptic} derive the curvature identity and solve the filled
system. Section~\ref{sec:height} locates the minimizing enclosure.
Sections~\ref{sec:weak3} and~\ref{sec:weak4} give the complete weak-decay
preparations and their distinct designated-end transfers.
The companion \emph{Boundary graph deformations for the spacetime
Penrose inequality} \cite{BoundaryCompanion} studies what changes when
one keeps an actual inner boundary, including its signed flux, or
retains a decaying tensor in a maximal-data construction. It provides
alternative analytic tools. The companion \emph{Equality and rigidity
in the spacetime Penrose inequality} \cite{RigidityCompanion} returns
to the original $g$ and $K$ through variation, a causal adjoint system
and static classification. These alternative methods and equality
results are not premises of the numerical argument proved here.

